\input{estilo}
\begin{document}
\begin{figura}\begin{tikzpicture}[srv/.style={caja, minimum width=8mm, minimum height=7mm},
    lee/.style={msg, draw=segundo}]
  % tres filas: el master (centro en y=2.75), los readers (y=0.35) y el load balancer (y=-1.9),
  % todas centradas en el reader del medio
  \pic at (2.5,2.2) {base};
  \foreach \x in {1, 2.5, 4} {
    \pic at (\x,-0.2) {base};
    \draw[msg] (2.5,2.05) -- (\x,1.05);
  }
  \node[caja] (lb) at (2.5,-1.9) {};
  \node[etiq, below=2pt] at (lb.south) {load balancer};
  \foreach \x in {1, 2.5, 4} \draw[lee] (lb.north) -- (\x,-0.4);
  % los servidores HTTP: el primero a la altura del master, el último a la del balancer
  \foreach \i/\y in {1/2.75, 2/1.2, 3/-0.35, 4/-1.9}
    \node[srv] (s\i) at (-1.6,\y) {};
  \node[etiq, align=center, below=2pt] at (s4.south) {servidores\\HTTP};
  \draw[msg] (s1.east) -- node[etiq, above] {escrituras} (2.0,2.75);
  \draw[lee] (s4.east) -- node[etiq, above, text=segundo] {lecturas} (lb.west);
  % el master copia cada escritura a sus readers
  \node[etiq, anchor=west] at (3.35,1.7) {replicación};
  % los nombres, cada uno a la altura de su fila
  \node[font=\small, anchor=west] at (5.2,2.75) {master / primary};
  \node[font=\small, anchor=west] at (5.2,0.35) {readers / backups};
\end{tikzpicture}\end{figura}
\end{document}
